\section{R ile çözümlü uygulama}
\label{sec:r-b06}

\textbf{Soru.} Her sınıfta dört öğrenci bulunan 12 kişilik çerçeveden altı
kişilik basit rastgele örneklem ve her sınıftan ikişer kişilik tabakalı
örneklem seçin. Seçilme olasılıklarını ve ağırlıkları bulun.

\begin{lstlisting}[style=prismR]
set.seed(2026)
ogrenciler <- data.frame(
  id = 1:12, sinif = factor(rep(1:3, each = 4))
)
basit <- ogrenciler[sample.int(nrow(ogrenciler), 6), ]
indisler <- split(seq_len(nrow(ogrenciler)),
                  ogrenciler$sinif)
secilen <- unlist(lapply(indisler,
                        function(indis) sample(indis, 2)),
                  use.names = FALSE)
tabakali <- ogrenciler[secilen, ]
tabakali$secim_olasiligi <- 2 / 4
tabakali$agirlik <- 1 / tabakali$secim_olasiligi
print(basit)
print(tabakali)
table(tabakali$sinif)
sum(tabakali$agirlik)
stopifnot(nrow(basit) == 6,
          anyDuplicated(basit$id) == 0,
          all(table(tabakali$sinif) == 2))
\end{lstlisting}

\begin{outputguidebox}
\textbf{Çözüm ve kontrol değerleri.} İki tasarımda da altı farklı öğrenci
seçilir. Tabakalı seçimde sınıf frekansları kesin olarak $(2,2,2)$'dir;
basit rastgele seçim bunu garanti etmez. Her öğrencinin seçilme olasılığı
$6/12=2/4=0{,}5$, tasarım ağırlığı 2'dir. Altı ağırlığın toplamı 12'dir.
\end{outputguidebox}

\textbf{Yorum.} Bu örnekte tabakalar eşit büyüklüktedir ve eşit sayıda
birim seçilir. Eşit olmayan tabaka hacimlerinde her tabaka için $N_h/n_h$
ayrı hesaplanmalıdır. Ağırlıklı ortalama hesaplamak, tasarıma uygun standart
hata hesabını kendiliğinden sağlamaz.

\textbf{Pekiştirme ve yanıt.} Bir sınıfın çerçeveden tamamen çıkarılması
kapsam hatası yaratır; kalan öğrencileri rastgele seçmek bu hatayı düzeltmez.